\section{Referee-round positive closure: typed parent spaces and proved tangent maps}
\label{sec:rr2-c2}

\subsection{A coproduct of physical platforms}

Let \(\mathfrak P\) be the finite set of platforms used in the series.  For
each \(p\in\mathfrak P\), let \((\Omega_p,\mathcal I_p)\) be its proved path
space and good rate.  Define the disjoint-union parent
\[
 \Omega^\sqcup=\bigsqcup_{p\in\mathfrak P}\{p\}\times\Omega_p,\qquad
 \mathcal I^\sqcup(p,\omega)=J(p)+\mathcal I_p(\omega),
\]
where \(J(p)\) is a fixed platform-selection cost.

\begin{theorem}[Typed coproduct contraction]
\label{thm:rr2-c2-coproduct}
Every platform-specific observable \(C_p:\Omega_p\to Y_p\) defines a
continuous typed map on its component.  Its contraction rate is
\[
 I_{p,C}(y)=\inf_{\omega:C_p(\omega)=y}\mathcal I_p(\omega).
\]
No rate from one platform is used on another.  If a genuine scaling map
\(S_n:\Omega_{p_n}\to\Omega_q\) is supplied, convergence of the pushed rates
is a separate theorem and is recorded explicitly.
\end{theorem}

\begin{proof}
The disjoint union topology makes each component open and closed.  The
ordinary contraction principle applied to the restriction gives the
formula.  The last assertion follows because a change of component is not a
continuous contraction unless a family of maps and laws is specified.
\end{proof}

\subsection{Precise cotangent quotient and duality}

For each platform choose the weighted Hölder path algebra \(\mathscr A_p\)
and the closed linear subspace
\[
 \mathscr N_p=
 \overline{\operatorname{span}}\left(
 \mathbb R\mathbf1+
 \{U-U\circ\Theta_t:U\in\mathscr A_p,\ t\in\mathbb Q\}
 \right).
\]
Equip \(\mathscr A_p\) with the weighted strict topology; its continuous dual
is the space of countably additive weighted Radon measures.

\begin{theorem}[Platform cotangent duality]
\label{thm:rr2-c2-duality}
The quotient \(\mathscr A_p/\mathscr N_p\) is Hausdorff.  For the pressure
\[
 Q_p(\Psi)=\sup_\nu\{\langle\Psi,\nu\rangle-\mathcal I_p(\nu)\},
\]
the subdifferential in the strict dual is exactly the compact set of
maximizing countably additive phases.
\end{theorem}

\begin{proof}
Closedness of \(\mathscr N_p\) gives a Hausdorff quotient.  Weighted strict
duality excludes finitely additive functionals.  Fenchel--Moreau duality and
compactness of rate sublevels identify subgradients with maximizing Radon
probabilities; conversely every maximizer satisfies the subgradient
inequality.
\end{proof}

\subsection{First and second pressure tangents}

\begin{theorem}[Spectral tangent theorem]
\label{thm:rr2-c2-tangent}
On a simple spectral phase of A2 or on a normally convergent cluster phase of
B2, \(Q_p\) is twice Fr\'echet differentiable on the declared source
algebra.  Its first derivative is the phase law, and
\[
 D^2Q_p(\Psi)[F,G]
 =
 \sum_{n\in\mathbb Z}
 \operatorname{Cov}_{\nu_\Psi}
 (F,G\circ\Theta_n)
\]
in collision time, with the corresponding Green--Kubo integral in physical
time.  The series/integral converges absolutely.
\end{theorem}

\begin{proof}
A2 supplies analytic dependence of the leading eigenvalue and eigenvectors;
B2 supplies normal convergence of an analytic cumulant series.  In either
case Cauchy estimates justify two differentiations.  Differentiating the
eigenvalue or cumulant expansion gives the correlation sum.  Exponential
mixing gives absolute convergence and identifies the physical-time integral
after the clock correction.
\end{proof}

\subsection{Rigidity with complete hypotheses}

\begin{theorem}[Dynamic entropic rigidity]
\label{thm:rr2-c2-rigidity}
On an atomless standard filtered probability space, let
\((\mathcal R_{s,t})\) be normalized, cash additive, relevant, law invariant,
continuous from above, and strongly time consistent.  If it is nonlinear,
then there is a constant \(\vartheta\ne0\) such that
\[
 \mathcal R_{s,t}(X)=
 \vartheta^{-1}\log
 \mathbb E[e^{\vartheta X}\mid\mathcal F_s].
\]
The theorem fixes the form but not the numerical coefficient.
On a mechanically calibrated cotangent work ray, unit covariance and
preparation consistency identify \(\vartheta\) with the mechanical field.
\end{theorem}

\begin{proof}
Cash additivity and law invariance reduce the one-step functional to a
law-based certainty equivalent.  Relevance and continuity rule out
degenerate affine pieces.  Strong time consistency gives the multiplicative
Cauchy equation for the certainty-equivalent transform, whose continuous
solutions are exponential or affine.  The nonlinear case is exponential
with a time-independent coefficient.  The final calibration follows from
the work-unit theorem on the chosen platform; it is not asserted for
unrelated payoffs.
\end{proof}

\subsection{Likelihood ratios, diffusion limits, and BSDEs}

For bounded path work \(F\),
\[
 M_t^\epsilon=
 \frac{\mathbb E[e^{\mu_\epsilon F}\mid\mathcal F_t]}
      {\mathbb E e^{\mu_\epsilon F}}
\]
is the microscopic likelihood-ratio martingale.

\begin{theorem}[Tangent likelihood-ratio convergence]
\label{thm:rr2-c2-likelihood}
Under the Gaussian tangent hypotheses of B3/D1, the pair consisting of the
fluctuation field and \(M^\epsilon\) converges jointly.  The limiting
likelihood ratio is
\[
 \mathcal E\!\left(
 \int_0^\cdot \vartheta\,\sigma^\top Du(s,X_s)\,dW_s
 \right),
\]
and is uniformly integrable.  The log-Laplace value solves the quadratic
BSDE
\[
 Y_t=\Phi(X_T)+\int_t^T
 \left(c_s+\frac{\vartheta}{2}|Z_s|^2\right)ds
 -\int_t^TZ_s\,dW_s .
\]
\end{theorem}

\begin{proof}
Uniform convergence of tilted cumulants on a complex neighborhood gives
mod-Gaussian convergence and exponential moment bounds.  Le Cam's third
lemma identifies the joint tilted Gaussian limit and its Cameron--Martin
shift.  The exponential moment bound implies Novikov and uniform
integrability.  Martingale representation of the Markov diffusion gives
\(Z=\sigma^\top Du\); applying It\^o's formula to the log-Laplace value
yields the BSDE, with uniqueness in the bounded-terminal quadratic class.
\end{proof}

\subsection{Memory compatibility}

\begin{proposition}[Commuting preparation diagram]
\label{prop:rr2-c2-memory}
On a fixed platform and phase, preparation commutes with: formation of the
history kernel, the resolvent memory Schur complement, and the terminal
likelihood-ratio tilt.  More precisely, all three are obtained by
disintegrating or compressing the same tilted path law, and their finite-time
conditional expectations form a commutative diagram.
\end{proposition}

\begin{proof}
Each arrow is a conditional expectation or a bounded compression under the
same Radon--Nikodym density.  The tower property and uniqueness of regular
conditional probabilities show that conditioning then tilting equals
tilting then conditioning.  The block resolvent identity is preserved under
the unitary identification of \(L^2\) spaces induced by the density.
\end{proof}
